\chapter{Totzeitmessung des GMZ}

\section{Aufbau}
\label{sec:gmz_theorie}
Für die weiteren Messungen, bei denen das Verhalten der Röntgenstrahlung bei Abschirmung betrachtet wird, wird ein Geiger-Müller-Zählrohr (GMZ) für den Nachweis der einzelnen Röntgenphotonen verwendet und hinten in den Experimentierraum des Röntgengeräts eingebaut. Ein GMZ ist ebenfalls ein gasgefüllter Ionisationsdetektor, der aus einer gasgefüllten Röhre (negativ geladene Außenwand) mit einem sehr dünnen Anodendraht in der Mitte besteht, welcher positiv geladen ist. Bei Inzidenz eines Röntgenphotons in das Gasvolumen werden dort Gasmoleküle ionisiert und aufgrund der $E(r)\propto\frac{1}{r}$-Abhängigkeit des elektrischen Felds in der Röhre (hier ist $r$ der Abstand zum Draht) stark zum Draht beschleunigt. Es kommt auf dem Weg aufgrund der kinetischen Energie durch das elektrische Feld zu Sekundärionisationen und damit Sekundärelektronen \cite[189-191]{kolanoskiwermerdetektoren}.

\jafps{ionisationsdet_bereiche}{Betriebsbereiche von gasgefüllten Teilchendetektoren, die mit Gasionisation funktionieren. Für das GMZ wird das angelegte elektrische Feld im Geiger-Bereich eingestellt. Für den Plattenkondensator in Aufgabenteil \todo{ref sec} wird der Ionisationskammerbereich genutzt. Abbildung entnommen aus \cite[194]{kolanoskiwermerdetektoren}.}{0.5}

Im Betriebsbereich (siehe die Betriebsbereiche in Abbildung \ref{fig:ionisationsdet_bereiche}) des Geiger-Müller-Zählrohrs ist die Feldstärke in der Nähe des Drahts sehr groß, wodurch dort für jedes einfallende Röntgenphoton immer eine maximal weit ausgebreitete (über die ganze Ausdehnung des Rohrs) und maximal starke Ladungslawine ausgelöst wird, die als großer Strompuls an der Anode registrierbar ist.
Jedes einfallende Photon wird so mit einem gleich hohen (maximal hohen) Stromimpuls registriert. Wenn über das ganze Zählrohr Gasmoleküle ionisiert sind, es also große Raumladungen gibt, neutralisieren diese das außen angelegte elektrische Feld und erlauben so die Rekombination von Elektronen-Loch-Paaren im Gas. Damit sinkt das Stromsignal am Anodendraht und das GMZ kann auf das nächste einfallende Röntgenphoton wieder mit neuen Ionisationen reagieren, also das nächste Signal messen. Aufgrund dieses Rekombinationsvorgangs dauert es eine endliche Zeit, bis ein GMZ nach einem Auslösen ein neues, getrenntes Signal registrieren kann, welche als \textit{Totzeit} bezeichnet wird. Bei Geiger-Müller-Zählrohren ist eine Totzeit nach einem einfallenden Teilchen von $\SIrange{50}{100}{\mu\s}$ zu erwarten, die maximale Teilchenrate, bei denen alle Teilchen korrekt registriert werden können, liegt damit bei $\SIrange{1e4}{1e5}{\hertz}$ \cite[194-196]{kolanoskiwermerdetektoren}.


\section{Durchführung}
Es wird mit dem GMZ die Teilchenrate bei verschiedenen Emissionströmen gemessen.

\section{Auswertung}
Es wird im Allgemeinen zwischen paralysierender und nicht paralysierender Totzeit unterschieden. Bei einer nicht paralysierenden Totzeit können einfallende Teilchen den unsensitiven Zeitraum nicht verlängern bevor dieser abgelaufen ist, die minimale Zeit zwischen zwei detektierten Ereignissen ist eine Konstante. Für die paralysierende Totzeit $\tau_p$ erwartet man zwischen gemessener Zählrate $m$ und tatsächlicher Zählrate $n$ den Zusammenhang \cite[776]{kolanoskiwermerdetektoren}
\begin{align}
	m &= \frac{n}{1 + n \tau_p}
\end{align}

Im Falle einer nicht paralysierenden Totzeit $\tau_n$ erwartet man den Zusammenhang\cite[776]{kolanoskiwermerdetektoren}
\begin{align}
	m = n \cdot e^{- n \tau_n}
\end{align}

Speziell für GMZ kann ein Mischfall betrachtet werden, bei dem eine paralysierende und eine nicht paralysierende Totzeit einfließen. \cite{gm_deadtime}
\begin{align}
	m = \frac{n \cdot e^{- n \tau_p}}{1 + \tau_n n}
\end{align}

Es wird hierbei angenommen, dass $n = k \cdot I_H$ gilt, da es nicht um die Dosis sondern nur die Anzahl der emittierten Photonen geht, welche proportional zum Heizstrom sein sollte.

Die Messdaten (in \tab{gmt_deadtime.table} zu finden) werden in \fig{deadtime} aufgetragen und die drei Totzeitmodelle mittels \textit{scipy curve\_fit} \cite{scipy} an die Daten angepasst. Betrachtet man in \fig{deadtime} die Anpassungsgüten, so erkennt man bei dem gemischten Totzeitmodell die beste Anpassungsgüte, bei dem paralysierenden Modell die schlechteste. Dies zeigt, dass die Totzeit des GMZ am besten durch das gemischte Modell beschrieben wird, als einfacheres Modell eine nicht paralysierende Totzeit allerdings auch relativ gut geeignet ist\footnote{zumindest bei den von uns gemessenen Teilchenraten}.

% \todo{fit bei nicht allen daten erläutern}

Die Fitergebnisse sind in \tab{deadtime_params} aufgelistet. Bei allen Modellen sind die Konstante $k$ in der gleichen Größenordnung, was Sinnvoll ist, da alle Modelle die gleiche Messung mit der gleichen Anzahl emittierter Teilchen beschreiben. Bei dem reinen nicht paralysierenden Modell findet man eine Totzeit $\tau_n$, wie sie für GMZ zu erwarten ist (vgl. \ref{sec:gmz_theorie}). Bei dem GMZ spezifsichen Modell findet man als $\tau_n$ eine Totzeit nahe der des einfachen nicht paralysierenden Modells, sowie eine sehr viel kürzere paralysierende Totzeit. Dies erscheint Sinnvoll, da in diesem Fall $\tau_p$ nur den Zeitraum abdeckt, in welchem im GMZ bereits wieder verstärkung möglich ist, das Signal aber nich nicht groß genug ist um von der Elektronik detektierbar zu sein.

\jafmt{|c|c|c|c|}{
	\hline
	Modell & $k$ / $\unit{\per\second\per\milli\ampere}$ & $\tau_p$ / $\unit{\second}$ & $\tau_n$ / $\unit{\second}$ \\
	\hline
	nicht paralysierend & \num{4.798(59)e+05} && $\num{8.794(28)e-05}$\\
	paralysierend &\num{1.2023(43)e+05} & \num{2.5625(57)e-05} & \\
	GMZ spezifisch & \num{3.931(54)e+05} & \num{6.51(34)e-07} & \num{7.601(55)e-05} \\\hline
}{Fit Parameter der verschiedenen Modelle.}{deadtime_params}

\jafps{deadtime}{Messung der Teilchenraten bei verschiedenen Emissionsströmen, mit Verschiedenen}{0.8}
% \jafps{deadtime_raw}{}{0.8}
% \jafps{deadtime_sliced}{}{0.8}
